% !TEX program = xelatex
\documentclass[12pt]{ctexart}
% ------------------------
% 页面
% ------------------------
\usepackage{geometry}
\geometry{margin=2.5cm}

% ------------------------
% 字体
% ------------------------
\usepackage{fontspec}
\setmainfont{Times New Roman}
\setmonofont{Consolas}

% ------------------------
% 基础包
% ------------------------
\usepackage{multirow}
\usepackage{soul}
\usepackage[normalem]{ulem}
\usepackage{lastpage}
\usepackage{xcolor}
\usepackage{graphicx}
\usepackage{booktabs}
\usepackage{enumitem}
\usepackage{caption}
\usepackage{titlesec}
\usepackage{fancyhdr}
\usepackage{hyperref}
\usepackage[backend=biber,style=gb7714-2015]{biblatex}
% ------------------------
% GitHub 颜色
% ------------------------
\definecolor{gh-bg}{HTML}{FFFFFF}
\definecolor{gh-border}{HTML}{D0D7DE}
\definecolor{gh-codebg}{HTML}{F6F8FA}
\definecolor{gh-link}{HTML}{0969DA}
\definecolor{gh-line}{HTML}{8B949E}

\definecolor{gh-note}{HTML}{0969DA}
\definecolor{gh-warning}{HTML}{9A6700}
\definecolor{gh-tip}{HTML}{1A7F37}

% ------------------------
% 设置页眉与页尾
% ------------------------
\geometry{left=2cm,right=2cm,top=2cm,bottom=2.5cm}
\fancyhf{}
\fancyhead[C]{}
\fancyhead[L]{中国科学院大学·2026年春}
\fancyhead[R]{网络空间安全导论}
\fancyfoot[C]{第 \thepage 页，共 \pageref{LastPage} 页}
\renewcommand{\footrulewidth}{0.4pt}
\renewcommand{\baselinestretch}{1.2}
\pagestyle{fancy}
% ------------------------
% 超链接
% ------------------------
\hypersetup{
	colorlinks=true,
	linkcolor=gh-link,
	urlcolor=gh-link,
	citecolor=gh-link,
	hyperfootnotes=true
}

% ------------------------
% 列表间距
% ------------------------
\setlist{
	itemsep=4pt,
	topsep=6pt
}

% ------------------------
% tcolorbox
% ------------------------
\usepackage{tcolorbox}
\tcbuselibrary{minted,skins,breakable}

% ------------------------
% GitHub 引用块
% ------------------------
\newtcolorbox{ghquote}{
	enhanced,
	colback=white,
	colframe=gh-border,
	borderline west={3pt}{0pt}{gh-border},
	arc=2pt,
	left=10pt,
	right=10pt,
	top=6pt,
	bottom=6pt
}

% ------------------------
% GitHub 提示框
% ------------------------
\newtcolorbox{note}{
	enhanced,
	colback=white,
	colframe=gh-note,
	borderline west={4pt}{0pt}{gh-note},
	title=Note
}

\newtcolorbox{warning}{
	enhanced,
	colback=white,
	colframe=gh-warning,
	borderline west={4pt}{0pt}{gh-warning},
	title=Warning
}

\newtcolorbox{tip}{
	enhanced,
	colback=white,
	colframe=gh-tip,
	borderline west={4pt}{0pt}{gh-tip},
	title=Tip
}

% ------------------------
% 代码块
% ------------------------
\usepackage{minted}

\newtcblisting{codeblock}[2][]{
	listing engine=minted,
	minted language=#2,
	minted options={
		linenos,
		breaklines,
		fontsize=\small,
		numbersep=10pt,
		xleftmargin=15pt
	},
	colback=gh-codebg,
	colframe=gh-border,
	arc=6pt,
	boxrule=0.5pt,
	left=6pt,
	right=6pt,
	top=6pt,
	bottom=6pt,
	listing only,
	breakable,
	#1
}

% ------------------------
% 行内代码
% ------------------------
\newtcbox{\inlinecode}{
	on line,
	boxrule=0pt,
	colback=gh-codebg,
	arc=4pt,
	left=3pt,
	right=3pt,
	top=2pt,
	bottom=2pt
}

% ------------------------
% 标题样式
% ------------------------
\titleformat{\section}
{\Large\bfseries}
{\thesection}{0.5em}{}

\titleformat{\subsection}
{\large\bfseries}
{\thesubsection}{0.5em}{}


\renewcommand{\theFancyVerbLine}{\color{gh-line}\small\arabic{FancyVerbLine}}

\newtheorem{application}{应用}
\newtheorem{problem}{问题}
\newtheorem{analogy}{分析}
\newtheorem{experiment}{实验}

% ------------------------
% 修改表头信息
% ------------------------
%改这里可以修改实验报告表头的信息
\newlength{\namegap} % 定义一个namegap的变量用来调节name之间的宽度，可以自由地修改这里来保证美观
\setlength{\namegap}{4.4em} 
\newcommand{\name}{
	fa\footnotemark[1] \hspace*{\namegap} 
	her\footnotemark[2] \hspace*{\namegap} 
	tms\footnotemark[3] \hspace*{\namegap} 
	tkz\footnotemark[4] \hspace*{\namegap} 
	aht\footnotemark[5] 
}
\newlength{\numgap} % 定义一个numgap的变量用来调节学号之间的宽度
\setlength{\numgap}{0.9em} 
\newcommand{\studentNum}{
	\scriptsize 
	\textbf{****K80099*****} \hspace*{\numgap}
	\textbf{****K80099*****} \hspace*{\numgap} 
	\textbf{****K80099*****} \hspace*{\numgap}
	\textbf{****K80099*****} \hspace*{\numgap}
	\textbf{****K80099*****} 
}
\newcommand{\expNum}{ICS-groupwork1}
\newcommand{\expName}{探究人工智能技术对网络空间安全的影响}
\newcommand{\finishdate}{xxxx年xx月xx日}
%直接使用语句 \tabularcenter{  } 针对性地对表格轴某一个格子居中
\newcommand{\tabularcenter}[1]{\multicolumn{1}{c}{#1}}
% ------------------------
% 文档开始
% ------------------------
\addbibresource{template_reference.bib}
%上面是调用外置的bib文件来执行引用
\begin{document}
	
	\begin{center}
		\LARGE 中国科学院大学·网络空间安全导论xx报告
	\end{center}
	
	\begin{center}
		\emph{实验序号} \underline{\makebox[9.5em][c]{\expNum}}\hspace*{10em}
		\emph{最后修改日期} \underline{\makebox[10em][c]{\finishdate}}\\
		\emph{实验专题} \underline{\makebox[36em][c]{\expName}}\\
		\makebox[4em][l]{\emph{小组成员\footnotemark[0]}} \underline{\makebox[36em][c]{\name}}\\
		\makebox[4em][l]{\emph{对应学号}} \underline{\makebox[36em][c]{\studentNum}}
		{\noindent}
		\rule[8pt]{17cm}{0.2em}
	\end{center}
	\footnotetext[0]{作者排名不分先后，按姓名拼音排序。}
	\footnotetext[1]{fa，学号(****K80099*****)，贡献：}
	\footnotetext[2]{her，学号(****K80099*****)，贡献：}
	\footnotetext[3]{tms，学号(****K80099*****)，贡献：}
	\footnotetext[4]{tkz，学号(****K80099*****)，贡献：}
	\footnotetext[5]{aht，学号(****K80099*****)，贡献：}
	\begin{abstract}
		这是一个摘要
	\end{abstract}
	\begin{warning}
		注意这个文件的命名格式是： 名字列表\_第<n>次大作业\_<实验标题>.pdf\\
		但是使用biber对中文的文件名不支持，所以推荐最后再手动处理或脚本(Renamer.py/.bat)处理，一开始就统一命名成homework<n>.tex的形式
	\end{warning}
	\section{这是第一个大Section}
	\subsection{这是第一个大Section下的第一个小section}
	\subsubsection{下面还有一个更小的section}
	\section{这是第二个Section}
	\section{展示引用}
	网络空间安全\cite{icsbook}
	\section{Code Example}
	
	\begin{codeblock}{bash}
git clone https://github.com/user/project
cd project
make
	\end{codeblock}
	下面是C++代码的展示
	\begin{codeblock}{C++}
#include<iostream>
using namespace std;
long long map[42][42] = { };
int horse_x,horse_y;
int dest_x,dest_y;
int iarray[4] = {-2 , -1 , 1 , 2};
int abx(int x){    return x>=0 ? x : -x;}
void mark(){
    map[0][0]=1;
    for(int i=0;i<4;i++){
        for(int j=0;j<4;j++){
            if(abx(iarray[i])!=abx(iarray[j]) && horse_y+iarray[j] >=0 && horse_x+iarray[i]>=0){
                map[horse_x+iarray[i]][horse_y+iarray[j]]=-1;
            }
        }
    }
    map[horse_x][horse_y] = -1;
}
int main(){
    cin>>dest_x>>dest_y>>horse_x>>horse_y;
    mark();
    for(int i=0;i<dest_x+dest_y+1;i++){
        for(int j=0;j<=i;j++){
            if(map[j][i-j]!=-1){
                if(map[j][i-j+1]!=-1)
                    map[j][i-j+1] += map[j][i-j];
                if(map[j+1][i-j]!=-1)
                    map[j+1][i-j] += map[j][i-j];
            }
        }
    }
    cout<<map[dest_x][dest_y];
    return 0;
}
	\end{codeblock}
	\twocolumn %表示从这里开始之后都是两栏格式的
	\section{Quote}
	
	\begin{ghquote}
		This is a style quote block.
	\end{ghquote}
	
	\section{Admonitions}
	
	\begin{note}
		This is a style note.
	\end{note}
	
	\begin{warning}
		Be careful with this command.
	\end{warning}
	
	\begin{tip}
		Use virtual environments.
	\end{tip}
	
	\section{Table}
	
	\begin{center}
		\begin{tabular}{lll}
			\toprule
			 \tabularcenter{Name} & \tabularcenter{Language} & \tabularcenter{Stars} \\
			\midrule
			Project C & BrainFuck(.,/>) & 10000 \\
			Project A & C++ & 1000 \\
			Project B & Python & 850 \\
			\bottomrule
		\end{tabular}
	\end{center}
	
	\section{List}
	
	\begin{itemize}
		\item Item one
		\item Item two
		\item Item three
	\end{itemize}
	%下面是调用外置的bib文件来执行引用
	\printbibliography
\end{document}
